\section{汇编器与 ELF：把未知量保留为显式证据}
\label{sec:assembler-elf}

汇编器不是把助记符逐字替换成操作码的字典。它解析节与符号，选择包括压缩形式在内的编码，计算局部布局，并把尚不能决定的地址表达为重定位（relocation）。由此得到的可重定位目标文件 \texttt{ET\_REL} 不是“半个可执行文件”，而是一份结构完整、专供链接器继续计算的契约。本节将概念模型与预检流程真实产生的 \texttt{handwritten-asm.rv64.o} 对照；所列数字来自随完整预检保存的 GNU \texttt{readelf}/\texttt{objdump} 报告，不把示意输出冒充实验结果。

\subsection{汇编的两遍模型：有用但不是实现限制}

经典的两遍汇编（two-pass assembly）模型很适合理解数据依赖。第一遍顺序扫描输入，维护位置计数器，为标签建立 \((\text{section},\text{offset})\) 符号值，并记录指示所建立的节；第二遍编码指令和数据，已知的同节距离可被求值，跨节或外部引用则产生重定位。真实汇编器可能合并扫描、迭代决定编码宽度，或采用片段（fragment）数据结构处理松弛；“两遍”是概念上的依赖分层，而非 GNU as 或 LLVM MC 必须恰好读文件两次。

这个模型解释了前向引用为何不需要源程序先声明。读到 \texttt{bge t1,a2,.Ldone} 时，\texttt{.Ldone} 尚未出现；等布局稳定后，汇编器知道它是同一 \texttt{.text} 中偏移 \texttt{0x38} 的局部符号。预检对象仍保留偏移 \texttt{0x4} 处的 \texttt{R\_RISCV\_BRANCH .Ldone}：这并不说明汇编器“算不出”，而是因为后续链接松弛可能删除前面的字节，重定位让链接器能在最终布局上重新编码分支。RISC-V psABI 明确要求参与松弛的序列携带相应重定位元数据\cite{bin-riscv-psabi-relax}。

\begin{figure}[t]
\centering
\begin{tabular}{ccl}
源汇编 & \(\longrightarrow\) & 节片段、符号、待求值表达式\\
 & 第一层提交 & 局部偏移、编码候选、对齐\\
 & \(\downarrow\) & \\
\texttt{ET\_REL} & = & 已编码字节 + 符号表 + 重定位表\\
\end{tabular}
\caption{汇编器的信息提交边界：未知地址不是填零后遗忘，而是以带类型的表达式延期。}
\label{fig:assembler-commit}
\end{figure}

\subsection{实物检查：对象头与节表}

\observed 对 \texttt{handwritten-asm.rv64.o} 的文件头检查得到：ELF64、小端、\texttt{EM\_RISCV}、\texttt{ET\_REL}，入口值为 0，无程序头表；标志为 \texttt{0x5}，即 RVC 与 double-float ABI。对象共有 10 个节，\texttt{.text} 大小 198 B、对齐 2 B；\texttt{.rela.text} 大小 432 B，恰为 18 个 \texttt{Elf64\_Rela} 项（每项 24 B）；\texttt{.symtab} 有 20 个 24 B 的符号项。没有程序头不是损坏：普通可重定位对象并不承担建立进程映像的职责。

\begin{table}[t]
\caption{\texttt{handwritten-asm.rv64.o} 中与论证有关的实测节}
\label{tab:objsections}
\centering
\small
\begin{tabular}{lrrp{.43\linewidth}}
\toprule
节 & 大小/B & 对齐/B & 作用\\
\midrule
\texttt{.text} & 198 & 2 & 指令字节；2 字节对齐体现 C 扩展可用\\
\texttt{.rela.text} & 432 & 8 & 18 个 RELA 重定位，\texttt{sh\_info=1} 指向被修补的 \texttt{.text}\\
\texttt{.data} / \texttt{.bss} & 0 / 0 & 1 & 案例没有静态存储期数据\\
\texttt{.note.GNU-stack} & 0 & 1 & 声明工具链栈属性的输入标记\\
\texttt{.riscv.attributes} & 69 & 1 & 架构扩展与兼容属性\\
\texttt{.symtab} & 480 & 8 & 20 个局部/全局/未定义符号项\\
\texttt{.strtab}, \texttt{.shstrtab} & --- & 1 & 符号名和节名字符串\\
\bottomrule
\end{tabular}
\end{table}

ELF 同一容器提供两个互相关联但目的不同的视图。节（section）是链接编辑视图，每个节头记录类型、标志、文件偏移、大小、对齐以及与别的节的关联；段（segment）是加载视图，由程序头描述一批将被共同映射的字节及其 R/W/X 权限\cite{bin-sysv-gabi}。\texttt{ET\_REL} 主要只有节；可执行文件或共享对象才需要用 \texttt{PT\_LOAD} 等段建立进程。一个段可覆盖多个节，节头甚至可从可运行的最终文件中剥离，而必要的程序头不能。因此“操作系统逐个加载 \texttt{.text}、\texttt{.data} 节”是方便但错误的缩写。

常见节也不能仅凭名字猜行为。\texttt{.bss} 通常是 \texttt{SHT\_NOBITS}：其内存大小可以非零，却没有对应文件字节；\texttt{.symtab}/\texttt{.strtab} 服务静态链接与调试，常在剥离后消失；\texttt{.dynsym}/\texttt{.dynstr} 才保存运行时动态解析所需的较小集合；\texttt{.debug\_*} 通常不带 \texttt{SHF\_ALLOC}，不会成为运行映像。节名是约定，真正的机器可读语义来自 \texttt{sh\_type}、\texttt{sh\_flags} 与链接字段。

\subsection{符号表：名字、定义位置和可见性是不同维度}

每个 \texttt{Elf64\_Sym} 含名字索引、值、大小、绑定、类型、可见性和定义节索引\cite{bin-sysv-gabi}。\texttt{STB\_LOCAL} 只在本对象内参与解析；\texttt{STB\_GLOBAL} 可跨对象匹配；\texttt{STB\_WEAK} 可被同名强定义覆盖，未解析弱引用通常取零。类型 \texttt{STT\_FUNC}/\texttt{STT\_OBJECT} 与绑定是正交维度，“全局函数”是二者的组合而非一种单独类别。\texttt{SHN\_UNDEF} 表示本对象提出需求，并不等价于编译错误。

\observed 案例对象中，\texttt{clipped\_dot} 是 \texttt{GLOBAL FUNC}，值 \texttt{0x0}、大小 60 B；\texttt{main} 同为 \texttt{GLOBAL FUNC}，值 \texttt{0x3c}、大小 138 B。\texttt{.Lloop}、\texttt{.Ldone} 等标签是定义在 \texttt{.text} 的局部符号。\texttt{getint}、\texttt{putint}、\texttt{putch} 则为全局 \texttt{SHN\_UNDEF}：对象说明“我需要这三个名字”，静态库稍后提供定义。反过来，\texttt{clipped\_dot} 即使在同一对象内已有定义，调用处仍可带重定位，以允许布局与松弛统一由链接器完成。

源文件的 \texttt{.type}/\texttt{.size} 因而并非装饰。如果漏掉它们，机器仍可能执行，但工具无法可靠显示函数边界，链接器与分析器也失去类型信息。类似地，符号值在 \texttt{ET\_REL} 中通常是“定义节内偏移”，不是将来的虚拟地址；把 \texttt{main=0x3c} 解释为进程将在地址 60 运行，就是把链接前后两个坐标系混为一谈。

\subsection{重定位是带类型的延期表达式}

RISC-V ELF64 使用显式加数的 \texttt{RELA}。每项记录应用位置 \(r_{offset}\)、符号和重定位类型，以及加数 \(A\)。概念上，链接器依据类型计算 \(S+A-P\)、\(S+A\) 或 \(B+A\)：\(S\) 是符号地址，\(P\) 是应用位置，\(B\) 是装载基址。但这不是把一个连续整数机械写入洞中。RISC-V 的 B/J/I/U 类立即数散布在指令位域，PC 相对地址还常拆成高 20 位与低 12 位；具体类型规定范围检查、舍入和位域编码\cite{bin-riscv-psabi}。

\observed 对象偏移 \texttt{0x80} 的源语句 \texttt{call getint} 被编码为

\begin{verbatim}
00000080: auipc ra, 0
           R_RISCV_CALL_PLT getint
           R_RISCV_RELAX
00000084: jalr  ra
\end{verbatim}

\texttt{CALL\_PLT} 表达一个可经 PLT 或直接定义完成的 PC 相对调用；伴随的 \texttt{RELAX} 授权链接器在保持语义时改写序列。\texttt{putint} 和 \texttt{putch} 的调用在 \texttt{0xac}/\texttt{0xb6} 呈现相同结构。局部循环回边则因已选用压缩跳转而具有 \texttt{R\_RISCV\_RVC\_JUMP}。因此对象中“立即数暂为零”不是缺失信息：完整答案是“指令占位字节 + 重定位类型 + 符号 + 加数”。只看 \texttt{objdump -d} 会漏掉后半份答案，应在对象阶段使用 \texttt{objdump -dr}。

RISC-V 的配对重定位尤其要求这种语义化表示。对于地址形成，\texttt{AUIPC} 所在的高位重定位与随后低位重定位共同表达 PC 相对地址；低位部分必须关联正确的高位位置。链接器还可能删除对齐 NOP 或把两条调用缩成一条。由此可见，重定位不是“把地址填进去”的无类型补丁，而是链接器可验证、可优化的代数承诺。

\subsection{静态归档：按需搜索的对象集合}

静态库 \texttt{.a} 是 ar 归档（archive），本质是若干成员文件与可选符号索引，不是预先合并好的单一 ELF。\texttt{ar t} 显示课程所给 \texttt{libsysy\_riscv.a} 只有成员 \texttt{sylib.o}；预检从 \texttt{lib/sylib.c} 为当前 Linux 交叉工具链重建的 \texttt{libsysy-glibc-rv64.a} 同样只有 \texttt{sylib.glibc.rv64.o}。\texttt{nm -s} 可查看索引，链接器借此直接找到定义所需符号的成员\cite{bin-gnu-ar}。

归档提取是需求驱动的。可用下面的简化状态机理解 GNU 链接器的传统扫描规则\cite{bin-gnu-ld}：

\begin{verbatim}
D := 已选择的全局定义；U := 尚未满足的强引用
从左到右扫描命令行：
  普通 .o：整体加入，更新 D 与 U
  归档 .a：反复抽取能满足当前 U 的成员，直到不再变化
结束：报告 U 中仍未定义的强符号及冲突强定义
\end{verbatim}

所以 \texttt{handwritten-asm.rv64.o} 应出现在运行时库之前：它先把 \texttt{getint}、\texttt{putint}、\texttt{putch} 放入 \(U\)，归档随后才有理由抽取 \texttt{sylib.glibc.rv64.o}。相反次序可能失败，因为扫描归档时尚无需求，链接器通常不会因后出现对象而自动回头。循环依赖可用 \texttt{--start-group}/\texttt{--end-group} 反复搜索，但有额外成本；\texttt{--whole-archive} 强制纳入全部成员，会增加代码与副作用，并非修复普通顺序错误的首选。

\subsection{四种工具回答四类问题}

表~\ref{tab:binarytools} 把实验命令与证据边界对应。工具输出不能替代解释：例如看到 \texttt{U putint}，正确结论是当前对象把定义延期给链接阶段，而不是“putint 不存在”；看到对象无程序头，也不能结论为“不是 ELF”。

\begin{table}[t]
\caption{二进制检查工具与可证实的问题}
\label{tab:binarytools}
\centering
\small
\begin{tabular}{lp{.34\linewidth}p{.39\linewidth}}
\toprule
工具 & 主要对象 & 本案例证据\\
\midrule
\texttt{readelf/readobj} & ELF 头、节、符号、重定位、程序头 & \texttt{ET\_REL}、10 节、18 个 \texttt{RELA} 项、RVC/LP64D 标志\\
\texttt{objdump -dr} & 指令解码并穿插重定位 & \texttt{call} 的 \texttt{AUIPC/JALR + CALL\_PLT/RELAX}\\
\texttt{nm} & 面向名字的定义/未定义摘要 & 对象中的 \texttt{U getint/putint/putch}；库中的 \texttt{T} 定义\\
\texttt{ar t}/\texttt{nm -s} & 归档成员与符号索引 & 运行时成员及其可满足的接口\\
\bottomrule
\end{tabular}
\end{table}

GNU \texttt{readelf} 有意不依赖 BFD，可作为调查工具缺陷时的第二解析器；LLVM 工具则与本实验的生成链自然配套\cite{bin-gnu-binutils,bin-llvm-tools}。最稳妥的复现记录应同时保存工具版本、完整命令和原始输出。下一节将沿 \(D/U\) 状态继续：解析只是链接器的第一步，此后还要做布局、重定位、松弛，并把面向链接的节重新组织为面向加载的段。
